% ============================================================================
%  SECCIÓN 5.3: SELECCIÓN DEL MECANISMO DE PERSISTENCIA
% ============================================================================

\section{Selecci\'on del mecanismo de persistencia}

Para garantizar que los datos cr\'iticos de la aplicaci\'on se conserven
m\'as all\'a del ciclo de vida de la sesi\'on, se evaluaron diferentes
mecanismos de persistencia disponibles en el ecosistema de React Native
y se seleccionaron los m\'as adecuados para cada tipo de dato.

\subsection{Evaluaci\'on de alternativas}

A continuaci\'on se presentan las alternativas evaluadas para la
persistencia en React Native:

\begin{table}[H]
  \centering
  \caption{Alternativas de persistencia evaluadas}
  \contenidoConFuente{%
    \begin{tablaAPA}
      \begin{tabular}{@{}p{2.6cm}p{2.8cm}p{2.2cm}p{2.0cm}p{4.8cm}@{}}
        \toprule
        \textbf{Mecanismo} & \textbf{Tipo} & \textbf{Duraci\'on} &
        \textbf{Complejidad} & \textbf{Uso en el proyecto} \\
        \midrule
        AsyncStorage & Key-value local & Persistente & Baja & Sesi\'on
        de usuario (tokens JWT). \\
        \addlinespace
        Supabase (PostgreSQL) & Base de datos relacional &
        Persistente & Media & Datos de dominio y favoritos. \\
        \addlinespace
        React Query & Cach\'e en memoria & Semi-permanente & Baja &
        Consultas y mutaciones de datos. \\
        \addlinespace
        useState & Estado local & Temporal & Muy baja & Formularios,
        filtros, UI. \\
        \addlinespace
        SQLite local & Base de datos local & Persistente & Alta &
        No utilizado (innecesario con backend). \\
        \addlinespace
        MMKV & Key-value local & Persistente & Media & No utilizado
        ( AsyncStorage suficiente). \\
        \bottomrule
      \end{tabular}
    \end{tablaAPA}%
  }{Elaboraci\'on propia.}
\end{table}

\subsection{Mecanismos seleccionados}

\subsubsection{AsyncStorage para la sesi\'on de usuario}

\textit{AsyncStorage} es el mecanismo de persistencia local recomendado
por React Native. Almacena datos en formato clave-valor de forma
as\'incrona y persiste entre sesiones. Se utiliza para almacenar el
token JWT y el refresh token de Supabase Auth:

\begin{lstlisting}
// Configuraci\'on del cliente Supabase
export const supabase = createClient(
  supabaseUrl, supabaseAnonKey, {
  auth: {
    storage: AsyncStorage,
    autoRefreshToken: true,
    persistSession: true,
  },
});
\end{lstlisting}

Las ventajas de AsyncStorage para este caso son:

\begin{itemize}
  \item Integraci\'on nativa con Supabase Auth.
  \item Persistencia autom\'atica de la sesi\'on.
  \item Soporte para renovaci\'on autom\'atica de tokens.
  \item Bajo consumo de recursos del dispositivo.
\end{itemize}

\subsubsection{Supabase (PostgreSQL) para datos de dominio}

Supabase proporciona una base de datos PostgreSQL gestionada con
una API REST autom\'atica (PostgREST) y pol\'iticas de seguridad
a nivel de fila (RLS). Se utiliza para almacenar todos los datos
de dominio de la aplicaci\'on:

\begin{lstlisting}
// Consulta de escenarios con relaciones
const { data, error } = await supabase
  .from('scenarios')
  .select(`
    *,
    scenario_images(url, is_primary),
    scenario_sports(sports(nombre))
  `)
  .eq('estado', 'activo')
  .order('nombre');
\end{lstlisting}

Las ventajas de Supabase para este caso son:

\begin{itemize}
  \item Base de datos relacional con soporte para consultas
    complejas y relaciones.
  \item Pol\'iticas RLS que garantizan la seguridad de los datos
    a nivel de fila.
  \item Tiempo real para actualizaciones instant\'aneas.
  \item Almacenamiento de archivos (im\'agenes) integrado.
  \item Autenticaci\'on gestionada con m\'ultiples m\'etodos.
\end{itemize}

\subsubsection{React Query para el cach\'e en memoria}

React Query gestiona el cach\'e de datos en memoria y sincroniza
autom\'aticamente el estado de la interfaz con el servidor. Se
utiliza para cachear las consultas a Supabase y gestionar las
mutaciones:

\begin{lstlisting}
// Hook con cach\'e de 5 minutos
export function useScenarios() {
  return useQuery({
    queryKey: ['scenarios'],
    queryFn: fetchScenarios,
    staleTime: 1000 * 60 * 5,
  });
}

// Invalidaci\'on tras mutaci\'on
await queryClient.invalidateQueries({
  queryKey: ['favorites', userId]
});
\end{lstlisting}

\subsection{Justificaci\'on de la combinaci\'on}

La combinaci\'on de AsyncStorage + Supabase + React Query cubre
los tres niveles de persistencia necesarios:

\begin{enumerate}
  \item \textbf{Persistencia local:} AsyncStorage conserva la
    sesi\'on del usuario entre aperturas de la aplicaci\'on,
    evitando que el usuario deba iniciar sesi\'on cada vez.

  \item \textbf{Persistencia en servidor:} Supabase almacena
    todos los datos de dominio de forma centralizada, permitiendo
    la sincronizaci\'on entre m\'ultiples dispositivos.

  \item \textbf{Cach\'e en memoria:} React Query reduce las
    consultas al servidor y mejora el rendimiento de la
    aplicaci\'on al cachear los datos m\'as utilizados.
\end{enumerate}
